% GENERATED FILE. Edit canonical lessons, then run npm run generate:books.
% canonical-source-hash: fnv1a64:b7fb35d6d0ed2b42
% canonical-lessons: PA-C342-samjhauna, PA-C342-matlab, PA-C342-karan, PA-C342-sapshat, PA-C342-sujhauna

\chapter{What Do You Mean?}
\label{ch:pa-a2-what-do-you-mean}
\hlchaptermodality{342}

\begin{chapteropening}
\textbf{By the end of this chapter:} \emph{I can ask what something means and give a reason.}

Complete the last lesson of chapter 342: I can ask what something means and give a reason.
\end{chapteropening}

\section[\texorpdfstring{\pa{ਸਮਝਾਉਣਾ} (samjhāuṇā)}{ਸਮਝਾਉਣਾ (samjhāuṇā)}]{\pa{ਸਮਝਾਉਣਾ} (samjhāuṇā) — to explain}

\label{lesson:PA-C342-samjhauna}



\begin{quote}
Before the new one: say the Punjabi for a guess, then the Punjabi for to seem right.
\end{quote}

\subsection*{You'll want to know: \pa{ਸਮਝਾਉਣਾ}}
\textbf{\pa{ਸਮਝਾਉਣਾ}} — \emph{samjhāuṇā} — \textquotedblleft{}to explain\textquotedblright{}.

\begin{quote}
\pa{ਕੀ} \pa{ਅਧਿਆਪਕ} \pa{ਹਰ} \pa{ਸਵਾਲ} \pa{ਸਮਝਾਉਂਦਾ} \pa{ਹੈ}?
\end{quote}

\begin{quote}
\emph{kī adhiāpak har savāl samjhāundā hai?}
\end{quote}

\begin{quote}
Does the teacher explain every question?
\end{quote}

Past (I …): \pa{ਮੈਂ} \pa{ਸਵਾਲ} \pa{ਸਮਝਾਇਆ}, \emph{maiṁ savāl samjhāiā}.

\begin{cousinweb}
\textbf{\pa{ਸਮਝਾਉਣਾ}} (\emph{samjhāuṇā}) is the causative partner of \emph{samajhṇā}, \textquotedblleft{}to understand\textquotedblright{}, which goes back to Sanskrit \emph{saṁbudhyate}, \textquotedblleft{}wakes up, understands\textquotedblright{}; CDIAL lists Punjabi \emph{samajhṇā} under it (CDIAL 12959).
\end{cousinweb}

\begin{grammarlens}[title={What you've built}]
One more word for asking and answering.
\end{grammarlens}

\subsection*{Your turn}
\begin{itemize}
\raggedright
  \item \emph{Say it:} \emph{samjhāuṇā}
  \item \emph{Say it:} \emph{samjhāuṇā}, once more
  \item \emph{Say it:} \emph{kī adhiāpak har savāl samjhāundā hai?} — Does the teacher explain every question?
  \item \emph{Recall:} say the Punjabi for a guess, then the Punjabi for to seem right, then say \emph{samjhāuṇā} again
\end{itemize}

\begin{tcolorbox}[breakable,colback=teal!4,colframe=teal!35!black,title={Before you move on}]
What does \pa{ਸਮਝਾਉਣਾ} mean? (\textquotedblleft{}To explain\textquotedblright{}.) Say it once more.
\end{tcolorbox}
\section[\texorpdfstring{\pa{ਮਤਲਬ} (matlab)}{ਮਤਲਬ (matlab)}]{\pa{ਮਤਲਬ} (matlab) — meaning}

\label{lesson:PA-C342-matlab}



\begin{quote}
Before the new one: say the Punjabi for to seem right, then the Punjabi for to explain.
\end{quote}

\subsection*{You'll want to know: \pa{ਮਤਲਬ}}
\textbf{\pa{ਮਤਲਬ}} — \emph{matlab} — \textquotedblleft{}meaning\textquotedblright{}.

\begin{quote}
\pa{ਕੀ} \pa{ਉਸ} \pa{ਸ਼ਬਦ} \pa{ਦਾ} \pa{ਕੋਈ} \pa{ਮਤਲਬ} \pa{ਹੈ}?
\end{quote}

\begin{quote}
\emph{kī us shabad dā koī matlab hai?}
\end{quote}

\begin{quote}
Does that word have any meaning?
\end{quote}

A masculine noun.

\begin{cousinweb}
\textbf{\pa{ਮਤਲਬ}} (\emph{matlab}) is a loan, usually traced through Persian to Arabic \emph{maṭlab}, \textquotedblleft{}something sought\textquotedblright{}.
\end{cousinweb}

\begin{grammarlens}[title={What you've built}]
One more word for asking and answering.
\end{grammarlens}

\subsection*{Your turn}
\begin{itemize}
\raggedright
  \item \emph{Say it:} \emph{matlab}
  \item \emph{Say it:} \emph{matlab}, once more
  \item \emph{Say it:} \emph{kī us shabad dā koī matlab hai?} — Does that word have any meaning?
  \item \emph{Recall:} say the Punjabi for to seem right, then the Punjabi for to explain, then say \emph{matlab} again
  \item \emph{Write it:} copy \textbf{\pa{ਮਤਲਬ}} once, checking the whole word
\end{itemize}

\begin{tcolorbox}[breakable,colback=teal!4,colframe=teal!35!black,title={Before you move on}]
What does \pa{ਮਤਲਬ} mean? (\textquotedblleft{}Meaning\textquotedblright{}.) Say it once more.
\end{tcolorbox}
\section[\texorpdfstring{\pa{ਕਾਰਨ} (kāran)}{ਕਾਰਨ (kāran)}]{\pa{ਕਾਰਨ} (kāran) — a reason; a cause}

\label{lesson:PA-C342-karan}



\begin{quote}
Before the new one: say the Punjabi for to explain, then the Punjabi for meaning.
\end{quote}

\subsection*{You'll want to know: \pa{ਕਾਰਨ}}
\textbf{\pa{ਕਾਰਨ}} — \emph{kāran} — \textquotedblleft{}a reason; a cause\textquotedblright{}.

\begin{quote}
\pa{ਕੀ} \pa{ਉਸ} \pa{ਦਾ} \pa{ਕੋਈ} \pa{ਕਾਰਨ} \pa{ਹੈ}?
\end{quote}

\begin{quote}
\emph{kī us dā koī kāran hai?}
\end{quote}

\begin{quote}
Is there a reason for that?
\end{quote}

A masculine noun.

\begin{cousinweb}
\textbf{\pa{ਕਾਰਨ}} (\emph{kāran}) is inherited from Sanskrit \emph{kāraṇa}, \textquotedblleft{}cause\textquotedblright{}; CDIAL lists Punjabi \emph{kāraṇ}, \textquotedblleft{}cause\textquotedblright{}, under it (CDIAL 3057).
\end{cousinweb}

\begin{grammarlens}[title={What you've built}]
One more word for asking and answering.
\end{grammarlens}

\subsection*{Your turn}
\begin{itemize}
\raggedright
  \item \emph{Say it:} \emph{kāran}
  \item \emph{Say it:} \emph{kāran}, once more
  \item \emph{Say it:} \emph{kī us dā koī kāran hai?} — Is there a reason for that?
  \item \emph{Recall:} say the Punjabi for to explain, then the Punjabi for meaning, then say \emph{kāran} again
  \item \emph{Write it:} copy \textbf{\pa{ਕਾਰਨ}} once, checking the whole word
\end{itemize}

\begin{tcolorbox}[breakable,colback=teal!4,colframe=teal!35!black,title={Before you move on}]
What does \pa{ਕਾਰਨ} mean? (\textquotedblleft{}A reason; a cause\textquotedblright{}.) Say it once more.
\end{tcolorbox}
\section[\texorpdfstring{\pa{ਸਪਸ਼ਟ} (sapshaṭ)}{ਸਪਸ਼ਟ (sapshaṭ)}]{\pa{ਸਪਸ਼ਟ} (sapshaṭ) — clear; plain to understand}

\label{lesson:PA-C342-sapshat}



\begin{quote}
Before the new one: say the Punjabi for meaning, then the Punjabi for a reason.
\end{quote}

\subsection*{You'll want to know: \pa{ਸਪਸ਼ਟ}}
\textbf{\pa{ਸਪਸ਼ਟ}} — \emph{sapshaṭ} — \textquotedblleft{}clear; plain to understand\textquotedblright{}.

\begin{quote}
\pa{ਤੁਹਾਡਾ} \pa{ਜਵਾਬ} \pa{ਸਪਸ਼ਟ} \pa{ਨਹੀਂ} \pa{ਹੈ।}
\end{quote}

\begin{quote}
\emph{tuhāḍā javāb sapshaṭ nahīṁ hai.}
\end{quote}

\begin{quote}
Your answer is not clear.
\end{quote}

\begin{cousinweb}
\textbf{\pa{ਸਪਸ਼ਟ}} (\emph{sapshaṭ}) is taken over from Sanskrit \emph{spaṣṭa}, \textquotedblleft{}clear\textquotedblright{}, a tatsama (Monier-Williams).
\end{cousinweb}

\begin{grammarlens}[title={What you've built}]
One more word for asking and answering.
\end{grammarlens}

\subsection*{Your turn}
\begin{itemize}
\raggedright
  \item \emph{Say it:} \emph{sapshaṭ}
  \item \emph{Say it:} \emph{sapshaṭ}, once more
  \item \emph{Say it:} \emph{tuhāḍā javāb sapshaṭ nahīṁ hai.} — Your answer is not clear.
  \item \emph{Recall:} say the Punjabi for meaning, then the Punjabi for a reason, then say \emph{sapshaṭ} again
  \item \emph{Write it:} copy \textbf{\pa{ਸਪਸ਼ਟ}} once, checking the whole word
\end{itemize}

\begin{tcolorbox}[breakable,colback=teal!4,colframe=teal!35!black,title={Before you move on}]
What does \pa{ਸਪਸ਼ਟ} mean? (\textquotedblleft{}Clear; plain to understand\textquotedblright{}.) Say it once more.
\end{tcolorbox}
\section[\texorpdfstring{\pa{ਸੁਝਾਉਣਾ} (sujhāuṇā)}{ਸੁਝਾਉਣਾ (sujhāuṇā)}]{\pa{ਸੁਝਾਉਣਾ} (sujhāuṇā) — to suggest}

\label{lesson:PA-C342-sujhauna}



\begin{quote}
Before the new one: say the Punjabi for a reason, then the Punjabi for plain to understand.
\end{quote}

\subsection*{You'll want to know: \pa{ਸੁਝਾਉਣਾ}}
\textbf{\pa{ਸੁਝਾਉਣਾ}} — \emph{sujhāuṇā} — \textquotedblleft{}to suggest\textquotedblright{}.

\begin{quote}
\pa{ਮੇਰਾ} \pa{ਦੋਸਤ} \pa{ਕਦੇ} \pa{ਕੋਈ} \pa{ਰਸਤਾ} \pa{ਨਹੀਂ} \pa{ਸੁਝਾਉਂਦਾ} \pa{ਹੈ।}
\end{quote}

\begin{quote}
\emph{merā dost kade koī rastā nahīṁ sujhāundā hai.}
\end{quote}

\begin{quote}
My friend never suggests a way out.
\end{quote}

Past (I …): \pa{ਮੈਂ} \pa{ਇੱਕ} \pa{ਰਸਤਾ} \pa{ਸੁਝਾਇਆ}, \emph{maiṁ ikk rastā sujhāiā}.

\begin{cousinweb}
\textbf{\pa{ਸੁਝਾਉਣਾ}} (\emph{sujhāuṇā}) is inherited from Sanskrit \emph{śudhyati}, \textquotedblleft{}is purified\textquotedblright{}, also \textquotedblleft{}becomes clear (of doubts)\textquotedblright{}; CDIAL lists Punjabi \emph{sujhāuṇā}, \textquotedblleft{}to explain\textquotedblright{}, under it (CDIAL 12527). CDIAL does not give the sense \textquotedblleft{}to suggest\textquotedblright{}.
\end{cousinweb}

\begin{grammarlens}[title={What you've built}]
One more word for asking and answering.
\end{grammarlens}

\subsection*{Your turn}
\begin{itemize}
\raggedright
  \item \emph{Say it:} \emph{sujhāuṇā}
  \item \emph{Say it:} \emph{sujhāuṇā}, once more
  \item \emph{Say it:} \emph{merā dost kade koī rastā nahīṁ sujhāundā hai.} — My friend never suggests a way out.
  \item \emph{Recall:} say the Punjabi for a reason, then the Punjabi for plain to understand, then say \emph{sujhāuṇā} again
  \item \emph{Write it:} copy \textbf{\pa{ਸੁਝਾਉਣਾ}} once, checking the whole word
\end{itemize}

\begin{tcolorbox}[breakable,colback=teal!4,colframe=teal!35!black,title={Before you move on}]
What does \pa{ਸੁਝਾਉਣਾ} mean? (\textquotedblleft{}To suggest\textquotedblright{}.) Say it once more.
\end{tcolorbox}
